\chapter{Perguntas}\label{cap:proposta}

As perguntas a serem respondidas no sistema final desenvolvido são:


\section{Quanto custa uma cadeira?}

Calcular a receita de cada candidato durante a campanha eleitoral e reter a média para cada cargo.

% \lipsum[67]

\section{Qual a taxa de sucesso eleitoral por patrimônio declarado?}

De acordo com o patrimônio declarado de cada candidato, estabelecer faixas arbitrárias de patrimônio e calcular a taxa de sucesso
 eleitoral para cada uma. Pode-se considerar dividir as faixas por quartis.

 \section{Quais os índices de dado município com relação aos eleitos/partidos mais vitoriosos?}

 Relacionar PIB per capita, IDHM, eleitores por população total, isentos (votos brancos + nulos + abstenções)
 aos partidos políticos mais votados. Pode-se calcular relação de viés político comparado à evolução de cada métrica em
 um gráfico.

 \section{Escolaridades do eleitorado e do candidato são correlatas?}

Calcular se existe correlação entre escolaridade do eleitorado e do eleito. Equiparar também
se existe relação entre escolaridade e taxa de sucesso em eleição.

\section{Quantos candidatos de dado partido foram eleitos por votos de legenda?}

Estabelecer ranking de partidos em relação a quantidade de eleitos por votos de legenda.

\section{Qual a relação entre a idade do eleitorado e do eleito?}

Calcular relação entre idade do eleitorado e idade do eleito.

\section{Como o viés político de um município evolui com o tempo?}

Analisar viés político de cada município com base nos partidos eleitos, atribuindo
pesos aos cargos para o cálculo final. Organizar ranking de municípios de acordo
com fidelidade política.

\section{Qual a relação da taxa de sucesso de um candidato com recebimento de recursos públicos?}

Analisar se candidatos que recebem mais recursos públicos tendem a ter mais chances de serem eleitos.

\section{Onde o candidato investe em propaganda?}

Estabelecer despesa de propaganda de cada candidato.

\section{Qual o histórico eleitoral de cada político?}

Expor vitórias, derrotas e reeleições de cada político. Mapear também probabilidade de reeleição.
